\documentclass[11pt]{article}
\usepackage{amsmath,amssymb,amsthm}
\newcommand{\I}{\mathsf I}\newcommand{\HH}{\mathsf H}
\newcommand{\Var}{\operatorname{Var}}
\newcommand{\supp}{\operatorname{supp}}
\newcommand{\E}{\mathbb E}
\newtheorem{proposition}{Proposition}
\begin{document}
\title{A002: exact transformation costs and endpoint admission}
\date{12 October 2026}\maketitle
All functions are signs on finite uniform product cubes. Entropies use bits.
Write $P_f(S)=\widehat f(S)^2$, $\HH(f)=H(P_f)$,
$\I_i(f)=\Pr(i\in S)$, $\I(f)=\E|S|$, and
$\Var(f)=1-(\E f)^2$. Constants have zero entropy and influence.
We never divide by the influence of a constant. Let
$C_*=\sup_{\I(f)>0}\HH(f)/\I(f)$, allowing infinity.

\section{Elementary ledgers}
For $f$ on $n$ bits, tensoring $m\ge1$ independent copies gives
$T_m=\prod_{j=1}^m f(x^{(j)})$ on $mn$ bits. Its mean is
$(\E f)^m$, entropy $m\HH(f)$, and influence $m\I(f)$.
Every coordinate copy keeps its original influence. The spectrum is a product
because the blocks are disjoint; overlapping products have no such ledger.

Balancing by a fresh sign gives $B_m=zT_m$ on $mn+1$ bits,
$\E B_m=0$, $\HH(B_m)=m\HH(f)$,
$\I(B_m)=m\I(f)+1$, and $\I_z=1$.
Its spectrum is the tensor spectrum with $z$ appended to every set.
Thus balancing does not make relative coordinate influence small.
Dummy padding by $q\ge0$ coordinates changes dimension by $q$ and changes
none of mean, entropy, variance or influence; every dummy has influence zero.

Multiplication by a character $\chi_R$ relabels $S$ by $S\triangle R$.
Entropy and distances between supports are unchanged; influences become
$\I_i$ or $1-\I_i$, according as $i\notin R$ or $i\in R$.
The new mean is $\widehat f(R)$. Selecting
$R=\{i:\I_i(f)>1/2\}$ gives total influence
$J=\sum_i\min(\I_i,1-\I_i)\le\I(f)$.

Balanced substitution has an exact ledger. If $u$ is on $N$ bits and
$g_i$ are balanced signs on disjoint blocks of sizes $d_i$, then
$U=u(g_1,\ldots,g_N)$ has dimension $\sum_i d_i$ and
\[
 \E U=\E u,\quad \HH(U)=\HH(u)+\sum_i\I_i(u)\HH(g_i),
 \quad \I(U)=\sum_i\I_i(u)\I(g_i),
 \quad \I_{ij}(U)=\I_i(u)\I_j(g_i).
\]
Indeed the set of occupied blocks uniquely identifies the outer Fourier set:
each inner spectral set is nonempty. Conditional on that outer set, the inner
sets are independent. The chain rule and expected cardinality prove all formulas.
The physical inputs remain independent even when the resulting spectral
membership coordinates are correlated.

\section{Point activation, including an unbalanced core}
Let $h$ be on $n$ bits with mean $a$, let $s\ge1$, let $A_s$ be one-point
activation on $s$ fresh bits, and put $p=2^{-s}$, $d=1-a$,
$G=1+A_s(h-1)$ on $n+s$ bits. Its exact ledger is
\[
 \E G=1-pd,\quad \Var(G)=2pd-p^2d^2,\quad
 \I(G)=p\I(h)+spd.
\]
Each data influence is $p\I_i(h)$; each control influence is $pd$.
For an all-positive active vertex the coefficients are $p\widehat h(S)$
at $(S,T)$ for $S\ne\varnothing$ and every $T$, $-pd$ at
$(\varnothing,T)$ for $T\ne\varnothing$, and $1-pd$ at the empty set.
Other active vertices only change the corresponding signs.
With $\phi(v)=-v\log_2v$,
\[
\begin{split}
 \HH(G)={}&p[\HH(h)-\phi(a^2)]+2sp(1-a^2)\\
 &+(p-p^2)d^2(2s-2\log_2d)+\phi((1-pd)^2).
\end{split}
\]
Terms at zero have their continuous value. This follows by summing the three
coefficient groups; the influence formula follows either from these coefficients
or from the pivotal event at each control.
For balanced $h$ the formula reduces to
\[
 \HH(G)=p\HH(h)+2sp(2-p)+\phi((1-p)^2),\qquad
 \I(G)=p[\I(h)+s].
\]
In particular $\phi((1-p)^2)/p\to2\log_2 e$.

\section{Systematic coding and its auxiliary costs}
Given a binary $r\times N$ matrix $A$, put
\[
 F_0(x,z)=G(x_i\chi_{A_i}(z):1\le i\le N).
\]
There are $N+r$ independent physical bits. The Fourier map is the injective
linear map $S\mapsto(S,AS)$, so
\[
 \E F_0=\E G,\quad\HH(F_0)=\HH(G),\quad
 \I(F_0)=\I(G)+\sum_S|AS|P_G(S)\le\I(G)+r\Var(G).
\]
The original influences are unchanged. Auxiliary bit $j$ has influence
$\Pr((AS)_j=1)\le\Var(G)$. Its spectral membership is a deterministic
parity of the original spectral memberships; this is not a physical bias.
If $r=1$ and $A$ is the all-one row, then the support has only even total
cardinality and $\I(F_0)=\I(G)+\Pr(|S|\text{ odd})\le2\I(G)$.
Its aggregate one-coordinate spectral erasure is zero; $z=1$ recovers $G$.
Repeating each original coordinate in a separate parity pair instead gives
$2N$ inputs, entropy $\HH(G)$ and influence $2\I(G)$ exactly.

For fixed integer $d_0\ge1$ and sufficiently large $N$, choose
$r=3d_0\lceil\log_2(2N)\rceil$. A uniformly random binary matrix has
\[
\Pr(\exists S\ne0:|S|+|AS|\le d_0)
\le2^{-r}\sum_{j=1}^{d_0}\binom Nj
        \sum_{l=0}^{d_0-j}\binom rl
\le2^{-r}d_0(N+r)^{d_0}<1.
\]
The middle bound follows by grouping weight-$j$ subsets in the $N$ original
coordinates and weight-$l$ subsets in the $r$ auxiliary coordinates, then
bounding their count by $d_0(N+r)^{d_0}$. For large $N$, $r\le N$.
Therefore a deterministic suitable matrix exists. This is an existence proof,
not a claim that a search of random matrices was executed.

\section{Diffuse seeds and auxiliary smoothing}
Let $\operatorname{Maj}_3=(x_1+x_2+x_3-x_1x_2x_3)/2$.
It has entropy $2$, total influence $3/2$, each influence $1/2$ and mean zero.
Character-center an arbitrary positive-entropy seed to $u$ with influence $J>0$.
For $q\ge1$ put $b_0=\operatorname{Maj}_3\prod_{j=1}^q u_j$,
on $n_0=3+qn$ bits. Then
\[
 L=\I(b_0)=qJ+3/2>1,\quad H_0=\HH(b_0)=q\HH(f)+2,
 \quad \max_i\I_i(b_0)\le1/2.
\]
Let $b_d$ be the balanced $d$-level self-substitution of $b_0$,
with $b_1=b_0$. It has $n_0^d$ bits and
\[
 \I(b_d)=L^d,\quad
 \HH(b_d)=H_0\frac{L^d-1}{L-1},\quad
 \max_i\I_i(b_d)\le2^{-d}.
\]
Choose $q$ large and then $d$ large. The limiting ratio is
$H_0/(L-1)$, which tends to $\HH(f)/J\ge\HH(f)/\I(f)$ as $q\to\infty$.
Consequently one fixed balanced diffuse seed $b$ has any prescribed small
maximum influence and ratio at least the original ratio minus any prescribed
positive error. If $\HH(f)=0$, a deep majority composition suffices.

For odd $L_g\ge3$ and integer $t\ge1$, define
$g=\operatorname{Maj}_{L_g}(\chi_{[t]}(z^{(1)}),\ldots,
\chi_{[t]}(z^{(L_g)}))$ on $d_g=L_gt$ bits. Then
\[
 \E g=0,\quad H_g=\HH(g)=\HH(\operatorname{Maj}_{L_g}),\quad
 I_g=\I(g)=tL_g\beta_{L_g},\quad
 \I_j(g)=\beta_{L_g}=2^{-(L_g-1)}\binom{L_g-1}{(L_g-1)/2}\to0.
\]
The pivotal identity follows by requiring a tie among the other majority inputs.
For $L_g=2v+1$, the central-binomial expression equals
$\prod_{j=1}^v(1-1/(2j))$. Its logarithm is at most
$-\tfrac12\sum_{j=1}^v1/j$, which tends to minus infinity; thus it tends to zero.
Majority is odd, so its nonzero supports have odd cardinality.
Parity substitution preserves its spectrum magnitudes. Every nonempty support
of $g$ has size at least $t$, and two distinct supports have distance at least $2t$.

Activate $h=b^{\otimes m}$, using $s=\lceil2\log_2m\rceil$ and $m\ge2$,
and then code all $N=mn+s$ inputs with a short-distance code having $r=O_k(\log m)$.
Set $e_j=\Pr((AS)_j=1)$ and $Q=sp+\sum_{j=1}^r e_j$.
Replace exactly the $s$ controls and $r$ coded auxiliary variables by independent
copies of $g$. The final dimension is $D=mn+d_g(s+r)$ and its exact ledger is
\[
\begin{split}
 \HH(F)&=pm\HH(b)+2sp(2-p)+\phi((1-p)^2)+H_gQ,\\
 \I(F)&=pm\I(b)+I_gQ,\qquad 0\le Q\le p[s+r(2-p)],\\
 \E F&=1-p,\qquad\Var(F)=p(2-p).
\end{split}
\]
An unchanged data bit has influence $p\I_i(b)$; a gadget bit in a control block
has influence $p\beta_{L_g}$; a gadget bit in coded block $j$ has influence
$e_j\beta_{L_g}$. Thus the data relative influence is
$\I_i(b)/(2-p)$ and every auxiliary relative influence is at most $\beta_{L_g}$.
All auxiliary energy has been charged in $I_gQ$; no auxiliary bit is free.
For fixed $b,g,k$, $D=mn+O(\log m)$,
$\I(F)\to0$, $\E F\to1$, and
$\HH(F)/\I(F)\to\HH(b)/\I(b)$.

Choose $d_0=t=2k+1$. Distinct coded outer supports differ on more than $d_0$
blocks; their substitutions differ on at least that many physical coordinates.
Distinct terms belonging to one outer support differ inside a gadget by at least
$2t$. Therefore the final Fourier-support distance exceeds $2k$.
Choose $b$ and $L_g$ to give maximum relative influence at most $\epsilon$,
then choose $m$ to satisfy total influence and near-constancy bounds.
This proves the companion manuscript's simultaneous reduction independently.

\section{Physical restriction plateau}
If support distance is greater than $q$, deleting any set $B$ of at most $q$
coordinates is injective on the support. Otherwise two supports differ only in $B$.
For every physical assignment $\rho$, the restricted coefficient at $T$ is
$\sum_{S\setminus B=T}\widehat F(S)\chi_{S\cap B}(\rho)$; at most one term
is nonzero. Hence
\[
 \HH(F|\rho)=\HH(F),\quad
 \I(F|\rho)=\I(F)-\sum_{i\in B}\I_i(F).
\]
Remaining influences are unchanged. The mean is unchanged if the empty support
was initially present; that holds for the rare activations above. Distance drops
by at most $|B|$. Thus every $\le k$ restriction of the constructed $F$ stays in
the simultaneous class with distance $>k$, and its ratio can only increase.
Such a child cannot be constant: it has the same mean strictly between zero and one.
This is finite-depth stability, not closure under all restrictions.

If a fixed nonnegative function $s_*(D)=o(D)$ is prescribed, the same construction
works with activation length
$\ell_m=\lceil a_m+2\log_2m\rceil$,
$a_m=\max_{1\le D\le2mn}s_*(D)$.
Indeed $a_m=o(m)$: split its maximum into a bounded prefix and a tail bounded
by any prescribed small multiple of $D$. Consequently $\ell_m=o(m)$,
$r=O(\log m)$, and final $D_m\le2mn$ eventually.
The error terms are $o(p_m m)$, so the ratio is unchanged, while
$\I(F_m)\le p_m m^2\le2^{-a_m}\le2^{-s_*(D_m)}$ for large $m$.
No monotonicity of $s_*$ was assumed. An exponential requirement
$\I\le2^{-cD}$ is not reached by this argument.

\section{Restricted endpoints and full Fourier support}
For balanced $g$, the selector
$M_g=((1+g)/2)y_1+((1-g)/2)y_2$ adds two bits and satisfies
\[
 \E M_g=0,\quad\HH(M_g)=2+\HH(g)/2,\quad
 \I(M_g)=1+\I(g)/2.
\]
The data influences halve and both fresh influences equal $1/2$. The coefficients
are $1/2$ at $\{y_1\},\{y_2\}$ and $\pm\widehat g(S)/2$ at
$S\cup\{y_1\},S\cup\{y_2\}$. For $g=z f^{\otimes m}$ a bound
$\HH(M_g)\le C\I(M_g)$ gives $m\HH(f)+4\le C[m\I(f)+3]$.
Division by $m\to\infty$ recovers FEI with the same constant.
Similarly activating $z f^{\otimes m}$ with
$s=\lceil2\log_2(m+1)\rceil$ gives arbitrarily small influence and minority
probability while its ratio tends to $\HH(f)/\I(f)$.

``Full Fourier support'' means all $2^D$ coefficients are nonzero; it is
different from relevance of every physical coordinate. In positive dimension
full Fourier support has distance one, so cannot be combined with distance $>1$.
It nevertheless retains the optimal FEI constant even together with rarity.

To prove this, take the last elementary activated tensor $G_m$ on $D_m=mn+1+s$
bits and pad it with two dummy bits, giving $M_m=D_m+2$ bits.
Flip its value at one physical point $x_0$ and call the result $K_m$.
Put $\delta_m=2^{-M_m}$. Every coefficient changes by
$-2\delta_m G_m(x_0)\chi_S(x_0)$.
All original nonzero coefficients have magnitude at least $2^{-D_m}$,
whereas $2\delta_m=2^{-D_m-1}$. Thus they cannot cancel; all original zero
coefficients become nonzero. All physical coordinates are consequently relevant.
The mean changes by at most $2\delta_m$, each influence by at most
$2\delta_m$, and total influence by at most $2M_m\delta_m$.
The influence bound follows because only the edge through $x_0$ in that
direction changes pivotal status; its two directed endpoints have mass
$2\delta_m$.

Parseval gives amplitude-vector distance $2\sqrt{\delta_m}$.
By Cauchy--Schwarz, squared-spectrum total variation $\tau_m$ is at most
$2\sqrt{\delta_m}$. For $0\le\tau_m\le1-2^{-M_m}$,
\[
 |\HH(K_m)-\HH(G_m)|\le h_2(\tau_m)+\tau_m\log_2(2^{M_m}-1).
\]
For completeness, a maximal coupling of the two finite laws disagrees with
probability $\tau_m$. Given one variable and the disagreement indicator,
the other has at most $2^{M_m}-1$ possibilities on disagreement. Bounding its
conditional entropy proves the inequality in both directions.
For large $m$, the same bound can use $2\sqrt{\delta_m}$ in place of $\tau_m$
because the right side is increasing up to $1-2^{-M_m}$.
Since $M_m\ge mn$ and $p_m$ is of order $m^{-2}$, these entropy and influence
errors are $o(p_m m)$. Hence
\[
 \I(K_m)\to0,\quad\E K_m\to1,\quad
 \HH(K_m)/\I(K_m)\to\HH(f)/\I(f).
\]
The full-support/low-influence class therefore has supremum $C_*$.
Applying the same surgery to the smoothed simultaneous sequence, with a strict
slack in its numerical hypotheses, also preserves small relative influences:
$|\Var(K_m)-\Var(F_m)|\le4\delta_m$ and
$\I_i(K_m)\le\I_i(F_m)+2\delta_m$, while $\delta_m/p_m\to0$.
It destroys the exact separated-support and plateau properties.

Finally fix any actual Boolean $F_*(X,Y)$ and tensor it with an arbitrary seed
on fresh $Z$ blocks: $F_m=F_*(X,Y)f^{\otimes m}(Z)$.
The $X$ spectral marginal and maximum $X$ Fourier degree are unchanged, while
\[
 H(S_{Y,Z}\mid S_X)=H_{F_*}(S_Y\mid S_X)+m\HH(f),\quad
 \I(F_m)=\I(F_*)+m\I(f),\quad
 \E F_m=(\E F_*)(\E f)^m.
\]
Dimension is $|X|+|Y|+mn$. A uniform conditional-entropy bound by $C\I(F_m)$,
or by $C\I_{Y,Z}(F_m)$ plus a fixed additive $X$ term, implies full FEI(C)
after division by $m$. Low $X$ degree or one fixed realizable $X$ marginal
does not exclude arbitrary seed tensors. A new admission hypothesis must exclude
this actual construction by proof, rather than by a requested conclusion.
No line above proves that $C_*$ is finite.
\end{document}
